% ---------------------------------------------------------------------
\section{末态相互作用与幺正性}\label{sec:disall-fsi}

在单个割图中，被击夸克与靶残余之间可以交换任意多个胶子。它们不破坏因子化的原因是幺正性：结构函数只依赖于对全部末态求和的结果，即前向振幅的不连续部分，而前向振幅中携带光子动量的线不被钳夹。

\begin{lemma}[末态喷注的非钳夹性 \tagA]\label{lem:disall-fsi}
对末态的全部割图求和给出前向振幅的不连续部分，$\Imag T_a(\omega+\ii0)=2\pi F_a$（$\omega>1$）。在前向振幅中，携带大分量 $q^-$ 的线不被钳夹：
在靶共线钳夹面附近，经过这些线的圈动量的负分量围道可以离开它们的极点一段 $\order{Q^2/P^+}$ 的距离，所以流喷注以及附着在它上面的软线都属于硬子图 $H$。
由\cref{eq:disall-current-virtuality}，$(\xi P+q)^2=Q^2(\xi-x)/x$：对 $\xi-x=\order{1}$，末态系统的不变质量为 $\order{Q^2}$，流喷注并入 $H$；
只在卷积的端点 $\xi\to x$ 处它才成为真正的共线喷注，这一端点产生系数函数中的加号分布，也是阈值区的来源。（本引理只涉及部分子态。）
\uses{prop:method-fsi,def:disall-breit,prop:disall-two-routes,prop:dis-compton-analytic}
% errtype: E-S
\end{lemma}

\begin{proofsketch}
(1) 一个前向图的全部末态割线之和等于它的不连续部分；按\cref{ch:dis}的归一化，这就是 $\Imag T_a=2\pi F_a$（\cref{prop:dis-compton-analytic}）。
(2) 携带 $q$ 的线 $k+q$ 的分母作为 $k^-$ 的函数，极点在 $k^-\simeq-q^-=-Q^2/(2x_NP^+)$ 附近，与共线钳夹面 $k^-=\order{\lambda^2Q^2/P^+}$ 相距 $\order{Q^2/P^+}$；
非物理区中的定量形式即\cref{eq:disall-unphysical-offshell}（\cref{prop:disall-two-routes}(a)）。附着在流喷注上的软动量 $\ell$ 只通过 $\ell^+$ 耦合（$(p'+\ell)^2-p'^2\simeq2p'^-\ell^+$），
极点都在 $\ell^+$ 平面的同一侧，围道可以移开：这是\cref{prop:method-fsi}的 DIS 情形。(3) 当 $\xi\to x$ 时流喷注的虚度趋于零，额外的末态共线区域给出加号分布（\cref{sec:disall-nlo}），
在固定 $x<1$ 处对光滑的 PDF 可积。完整论证见\cite[第~8.7~节]{CollinsSoperSterman1989}与\cite{LibbySterman1978a}。
\end{proofsketch}

\paragraph{地位与缺口}
这是 E-S 型论证：它说明末态钳夹面不是领头的，但不给出余项的积分界。固定 $x$ 处所需的割图版本建立在\cref{asm:appN-cutsum}之上（\cref{prop:disall-two-routes}(c)）；
一般的相消机制\cref{prop:method-fsi}（全阶论证，[A]）还依赖光锥序表示的假设\cref{asm:appO-lcopt}。文献中这一步是标度陈述：\cite[第~8.7~节]{CollinsSoperSterman1989}写道
「“the sum over final states eliminates pinch singularities involving final state jets”」，即对末态求和消除了涉及末态喷注的钳夹奇点。以\cref{asm:appN-cutsum}为前提时，本引理的割图形式已包含在\cref{prop:disall-two-routes}(c) 的 $R_{\mathrm{cond}}$ 陈述之中；提升本引理的标签则需要对 DIS 的割图证明该假设。

\begin{proposition}[DIS 的领头区域 \tagA]\label{prop:disall-regions}
$T^{\mu\nu}$ 的领头区域（等价地，对割图求和后 $W^{\mu\nu}$ 的领头区域）由一个靶共线子图与一个硬子图组成：二者由两条部分子（两条夸克线，或两条横向极化的胶子）外加任意多条 $A^+$ 极化（K）胶子相连；
没有软子图；对割图求和后没有末态共线子图。（本命题只涉及部分子态；强子版本即\cref{asm:intro-bridge}(Br1)。）
\uses{thm:ir-leading-regions,lem:disall-fsi}
% errtype: E-S
\end{proposition}

\begin{proofsketch}
(1) 对每个割图应用\cref{thm:ir-leading-regions}（全阶论证，[A]）：领头区域由靶共线子图 $J_A$、流喷注 $J_q$ 与只以胶子附着在二者上的软子图组成，没有软线附着在硬子图上（\cref{fig:reduced-dis}）。
(2) 由\cref{lem:disall-fsi}，割图求和后 $J_q$ 与附着在其上的软线并入硬子图；附着在硬子图上的每条软胶子压低 $\lambda^2$，只连接 $J_A$ 自身的软线是共线子图内部的子区域。
(3) 活动部分子为胶子时，横向胶子与 K 胶子组合为 $k^+A^j-k^jA^+$，即胶子 PDF 算符中 $G^{+j}$ 的 Fourier 分量。文献：\cite{LibbySterman1978a}，\cite[第~8.7~节]{CollinsSoperSterman1989}。
\end{proofsketch}

\paragraph{地位与缺口}
这是标度陈述（E-S）：嵌套区域没有 Minkowski 空间中与 Weinberg 定理相当的结果，个别 Feynman 规范图中的超领头项只在对图求和后相消。
其积分形式是\cref{asm:appN-envelope}（幂次计数包络）与\cref{asm:appN-separation}（分离性），作为\cref{thm:disall-factorization}的前提出现；固定 $x$ 处的割图版本依赖于\cref{asm:appN-cutsum}。
文献中的幂次计数\cite{LibbySterman1978a}与\cite[第~5~章]{Collins2011}给出的是标度论证，从未在所有阶写成积分形式；解除途径是幂次计数的积分形式（逐区域的扇形分解）加上分离性所需的 Łojasiewicz 型不等式。
